\begin{exercise}{97}
In this exercise, you may use the fact that the unit map
$K\to\mathcal S(|K|)$ is a homotopy equivalence for any Kan complex $K$.
Recall that a map of simplicial sets is a weak equivalence if its geometric
realization is a homotopy equivalence, and let $X$ be a simplicial set.
Prove or disprove the following statements:
\begin{enumerate}[label=(\arabic*)]
\item The unit map $X\to\mathcal S(|X|)$ is a monomorphism.
\item The unit map $X\to\mathcal S(|X|)$ is a weak equivalence.
\item The unit map $X\to\mathcal S(|X|)$ is anodyne.
\end{enumerate}

\begin{proof}
All three statements are true. Write $\eta$ and $\varepsilon$ for the unit
and counit of $|{-}|\dashv\mathcal S$ (Proposition~1.1.36).

(1) The singular simplex $\eta_X(\sigma)$ is the characteristic map
$t\mapsto[\sigma,t]$.
Suppose $\eta_X(\sigma)=\eta_X(\tau)$ for $\sigma,\tau\in X_n$.
By Exercise~\ref{ex:5}, write $\sigma=\alpha^*x$ and $\tau=\beta^*y$,
where $x,y$ are nondegenerate and $\alpha,\beta$ are surjections.
For interior points $t$ of the topological $n$-simplex, their affine images
$\alpha_*(t)$ and $\beta_*(t)$ are also interior points.
The open cells of $|X|$ are the interiors of its nondegenerate simplices
(Corollary~1.1.41).
Their disjointness and the injectivity of each characteristic map on its
interior imply $x=y$ and $\alpha_*(t)=\beta_*(t)$ for every such $t$.
By continuity the affine maps agree on the whole simplex; evaluating at
its vertices gives $\alpha=\beta$.
Hence $\sigma=\tau$, proving injectivity in every degree.

(2) Put $R=\mathcal S|{-}|$, $Y=|X|$, $Z=|RX|$,
$r=|\eta_X|\colon Y\to Z$ and $c=\varepsilon_Y\colon Z\to Y$.
The complex $RX$ is Kan by Lemma~1.1.47.
The granted hypothesis therefore makes
$a=|\eta_{RX}|\colon Z\to|\mathcal S Z|$ a homotopy equivalence.
Here realization sends simplicial homotopies to topological homotopies,
using the prism identification
$|K\times\Delta^1|\cong|K|\times[0,1]$.
The triangle identities give
\[
 cr=\mathrm{id}_Y,\qquad
 |\mathcal S c|\,a=\mathrm{id}_Z=\varepsilon_Za.
\]
In the homotopy category of spaces, $a$ is invertible, so
$|\mathcal S c|=\varepsilon_Z$.
Naturality of the counit now gives, in that category,
\[
 rc=\varepsilon_Z|\mathcal S r|
   =|\mathcal S c|\,|\mathcal S r|
   =|\mathcal S(cr)|=\mathrm{id}_Z.
\]
Thus $c$ is a homotopy inverse to $r$, proving the weak equivalence.

(3) By Fact~1.3.17, a monomorphism whose realization is a weak equivalence
is anodyne. Apply this to (1) and (2).
\end{proof}
\end{exercise}
